By \(\mathrm{ass:positivity}\), \(0<e<1\), and hence \(c=1/e>1\).  Therefore
\[
f_c(c)=(c-c)_+=0,
\qquad
f_c(0)=(0-c)_+=0.
\]
Using the definition of \(B_f(e)\), it follows that
\[
B_{f_c}(e)=e f_c(1/e)+(1-e)f_c(0)
=e f_c(c)+(1-e)f_c(0)=0.
\]

Put the two points of the Bernoulli outcome space in the order \(1,0\).  Since
\(p,q\in(0,1)\), both \(P\) and \(Q\) assign strictly positive mass to both
points.  Thus \(P\sim Q\), as required by \(\mathrm{def:jk\text{-}ball}\), and
the two Radon--Nikodym ratios are
\[
\frac{dP}{dQ}(1)=\frac pq,
\qquad
\frac{dP}{dQ}(0)=\frac{1-p}{1-q}.
\]
The definition of \(D_f(P\|Q)\) consequently gives
\[
\begin{aligned}
D_{f_c}(P\|Q)
&=q\left(\frac pq-c\right)_+
 +(1-q)\left(\frac{1-p}{1-q}-c\right)_+\\
&=(p-cq)_+ + ((1-p)-c(1-q))_+.
\end{aligned}
\]
Here the second equality uses \(s(x/s-c)_+=(x-cs)_+\), first with
\(s=q>0\) and then with \(s=1-q>0\).

Both summands in the last display are nonnegative, so their sum is zero if and
only if each is zero.  Since \(x_+=0\) if and only if \(x\le 0\), and since
\(1/c=e\),
\[
\begin{aligned}
D_{f_c}(P\|Q)=0
&\Longleftrightarrow p-cq\le0
 \ \text{and}\ (1-p)-c(1-q)\le0\\
&\Longleftrightarrow q\ge \frac pc=ep
 \ \text{and}\ q\le 1-\frac{1-p}{c}=1-e+ep.
\end{aligned}
\]
It remains to verify that this interval is exactly the fiber specified by
\(\mathrm{def:mixture\text{-}class}\).  If
\(Q=eP+(1-e)R\) and \(R=\operatorname{Bernoulli}(r)\), then evaluation at
the point \(1\) yields
\[
q=ep+(1-e)r.
\]
Because \(r\in[0,1]\), this implies \(ep\le q\le ep+1-e\).  Conversely, if
\(ep\le q\le1-e+ep\), define
\[
r=\frac{q-ep}{1-e}.
\]
The denominator is strictly positive by \(\mathrm{ass:positivity}\), and the
two interval inequalities give \(0\le r\le1\).  Hence
\(R=\operatorname{Bernoulli}(r)\) is a probability law and
\(Q=eP+(1-e)R\).  Moreover, because \(p\in(0,1)\), \(P\) has full support on
the binary space, so every such \(R\) satisfies \(R\ll P\).  This proves the
claimed equality with the binary fiber of \(\mathfrak L_e(P)\).

Finally, let \(\lambda>0\) and \(b\in\mathbb R\), and define
\(\widetilde f(t)=\lambda f_c(t)+b(t-1)\).  Since \(c>1\), one has
\(f_c(1)=0\); positive scaling preserves convexity and addition of an affine
function preserves convexity, so \(\widetilde f\) is finite, continuous,
convex, and \(\widetilde f(1)=0\).  In the terminology of
\(\mathrm{def:hinge\text{-}class}\), subtracting \(b(t-1)\) leaves
\(\lambda f_c(t)\), which vanishes on \([0,c]\) and is strictly positive on
\((c,\infty)\).  Its radius is still zero because \(ec=1\):
\[
\begin{aligned}
B_{\widetilde f}(e)
&=e\{\lambda f_c(c)+b(c-1)\}
 +(1-e)\{\lambda f_c(0)-b\}\\
&=b\{e(c-1)-(1-e)\}=b(ec-1)=0.
\end{aligned}
\]
Also, since \(P\ll Q\) and both are probability laws,
\[
\begin{aligned}
D_{\widetilde f}(P\|Q)
&=\lambda D_{f_c}(P\|Q)
 +b\int_{\mathcal Y}\left(\frac{dP}{dQ}-1\right)dQ\\
&=\lambda D_{f_c}(P\|Q)+b\{P(\mathcal Y)-Q(\mathcal Y)\}
=\lambda D_{f_c}(P\|Q).
\end{aligned}
\]
As \(\lambda>0\), the last quantity vanishes exactly when
\(D_{f_c}(P\|Q)\) vanishes.  Thus positive scaling and addition of
\(b(t-1)\) leave both the zero radius and the exact binary legality fiber
unchanged.